\documentclass{situationist-preprint}

\title{Feyerabend's \enquote{Anything Goes} Reconsidered}
\msid{SCI-05-b}
\msfamily{SCI-05}
\msdiscipline{Philosophy of Science}
\mstype{Companion manuscript}
\mskeywords{Feyerabend; anything goes; scientific method; Galileo; methodological norms}

\begin{document}
\maketitle

\begin{abstract}
Feyerabend's slogan \enquote{anything goes} is usually read as the claim that science has no method. We argue that, read as he presented it, the slogan is the conclusion of a reductio against the demand for universal methodological rules, and that it is compatible with the conclusion of SCI-05: there is no single method, but a family of methods united by common norms. Feyerabend's central case, Galileo's defense of the moving earth, violated particular rules of method but satisfied the norms of answerability and public checkability, and where it fell short of the norm of proportion the norms explain why Galileo's contemporaries were right to withhold acceptance and right to pursue his theory.
\end{abstract}

\section{The slogan}\label{s:slogan}

\citet{Feyerabend1975} argued that every methodological rule proposed by philosophers, including the rules to accept only theories consistent with well-confirmed observations and to reject theories refuted by experiment, has been broken in episodes now counted as scientific advances, and that the advances would not have occurred had the rules been followed. He concluded that the only principle that does not inhibit progress is \enquote{anything goes}. He also said that the slogan was not his own methodology but the position into which a rationalist who insists on universal rules is driven by the history \citep{Feyerabend1978}.

So read, the argument has the form: if there is a method, it is a set of universal rules; every universal rule has legitimate exceptions; so there is no method, and nothing is excluded. SCI-05-a accepts the second premise for rules, which it calls methods (Definition~3.1 of SCI-05-a), and its Proposition~4.1 is that no single method is both contentful and general. It rejects the first premise. Norms, in the sense of Definition~3.2 of SCI-05-a, are standards every method must satisfy without specifying a procedure, and the history Feyerabend cites shows exceptions to rules, not to norms.

\section{Galileo}\label{s:galileo}

Galileo's answer to the tower argument, that a stone dropped from a tower would land far to the west if the earth rotated, introduced a new interpretation of the observation: the stone shares the earth's motion, and shared motion is not perceived. Feyerabend argued that Galileo won assent by rhetoric and propaganda, without independent evidence for the new interpretation, and in violation of the rule to retain theories consistent with observation.

\begin{proposition}\label{p:galileo}
Galileo's procedure violated particular rules of method, but satisfied Principles~P1 (answerability) and P2 (public checkability) of SCI-05-a; its shortfall with respect to P3 (proportion) concerned acceptance, not pursuit.
\end{proposition}

\begin{proof}
The heliocentric theory, with its new account of shared motion, was answerable: it implied the phases of Venus, which Galileo observed and which the Ptolemaic system excludes, and it was later exposed to tests that could have refuted it, such as stellar aberration and stellar parallax. It was publicly checkable: Galileo published his arguments in the vernacular and invited others to observe through the telescope. Where his confidence outran the evidence available in 1632, as Feyerabend rightly emphasizes, the norm of proportion records the fault; but proportion governs the confidence with which a theory is accepted, and a theory not yet warranted may be worth pursuing \citep{Laudan1977}.
\end{proof}

\begin{objection}[Redescription]
Any episode can be made to satisfy norms as abstract as answerability and checkability. The norms are as empty as \enquote{anything goes}.
\end{objection}

\begin{reply}
They are not empty, since by Proposition~4.3 of SCI-05-a the characteristic procedures of pseudo-science violate them: a doctrine that no evidence could count against fails answerability, and a claim supported only by a private revelation fails checkability. What Feyerabend shows is that the norms exclude less than the rules, not that they exclude nothing.
\end{reply}

\section{Conclusion}\label{s:conclusion}

\enquote{Anything goes} is the correct verdict on the search for universal rules and the wrong verdict on science. Feyerabend's cases are counterexamples to rules, and the norms of SCI-05 survive them.

\begin{thebibliography}{9}
\bibitem[Feyerabend(1975)]{Feyerabend1975} Feyerabend, P. (1975). \emph{Against Method: Outline of an Anarchistic Theory of Knowledge}. New Left Books.
\bibitem[Feyerabend(1978)]{Feyerabend1978} Feyerabend, P. (1978). \emph{Science in a Free Society}. New Left Books.
\bibitem[Laudan(1977)]{Laudan1977} Laudan, L. (1977). \emph{Progress and Its Problems: Towards a Theory of Scientific Growth}. University of California Press.
\end{thebibliography}

\end{document}
